\section{Proof of Theorem B: The ideals arising from numerical semigroup rings}
\label{sec:target-ideals}


In this section, we give a proof of Theorem B. 
Throughout the remainder of this paper, let
$H=\langle a_1,a_2,a_3\rangle$ be a numerical semigroup of embedding
dimension three, and retain the notation $S$, $\varphi_H$, and $\fkp_H$
introduced in Subsection~\ref{subsec:basicnotion}.  Put $R=k[H]$, and let $\fkm=(t^{a_1},t^{a_2},t^{a_3})R$ denote the irrelevant maximal ideal of $R$.
For $h\in\mathbb N$, define
$$
 I_h=\varphi_H^{-1}\bigl(t^hk[t]\cap R\bigr).
$$
Equivalently, $I_h$ is generated by $\fkp_H$ together with all monomials of
weighted degree at least $h$.
Let
\[
 \ell=\min\{\deg\rho\mid
 \rho\in\fkp_H\setminus\{0\}\text{ is homogeneous}\}.
\]

Although $I_{\ell+1}$ is defined as an inverse image, this description does
not control its powers.  Indeed,
$\varphi_H^{-1}((t^{\ell+1}k[t]\cap R)^s)=I_{\ell+1}^s+\fkp_H$, which need not
equal $I_{\ell+1}^s$.  In particular, a relation in $\fkp_H$ of degree
$\ell$ belongs to $\varphi_H^{-1}((t^{\ell+1}k[t]\cap R)^2)$ but not to
$I_{\ell+1}^2$, whose nonzero homogeneous elements have degree at least
$2\ell$.  Thus the normality of $I_{\ell+1}$ cannot be deduced merely by
taking inverse images of the powers of $t^{\ell+1}k[t]\cap R$.

\begin{proposition}
\label{prop:inverse-image-ideals}
The following assertions hold.
\begin{enumerate}
	\item The ideal $I_h$ is integrally closed.
	\item If $H$ is symmetric, then $\mu_S(I_h) \le \mult(H) + 2$.
	\item If $H$ is not symmetric, then $\mu_S(I_h) \le \mult(H) + 3$.
	\item $I_h$ is a monomial ideal if and only if $h\leq\ell$.
\end{enumerate}
\end{proposition}

\begin{proof}
(1) The ideal $t^hk[t]$ is integrally closed in $k[t]$.  Its contraction
$t^hk[t]\cap k[H]$ is therefore integrally closed in $k[H]$, and its inverse
image $I_h$ is integrally closed in $S$.


For (2) and (3), the case $h=0$ is immediate, since $I_0=S$. Assume that $h>0$, and put $J=t^hk[t]\cap R$. Then $J$ is an $\fkm$-primary ideal of $R$.  For each $0 \le i \le \mult(H)-1$, choose the least element $h_i$ of $H$ such that $h_i \equiv i \pmod{\mult(H)}$ and $h_i \ge h$.  Then $J=(t^{h_i} \mid 0 \le i \le \mult(H)-1)$. Hence $\mu_R(J)\leq\mult(H)$.
Choose generators $f_1,f_2,\ldots,f_q$ of $J$ with $q\leq\mult(H)$, and choose $g_1,g_2,\ldots,g_q\in S$ such that $\varphi_H(g_i)=f_i$ for every $1\leq i\leq q$.  Since $I_h=\varphi_H^{-1}(J)$, we have
\[
 I_h=\fkp_H+(g_1, g_2, \ldots , g_q).
\]
By Theorem~\ref{thm:herzog}, $\fkp_H$ can be generated by two elements if $H$ is symmetric and by three elements otherwise.  Consequently, $\mu_S(I_h)\leq\mult(H)+2$ in the symmetric case and $\mu_S(I_h)\leq\mult(H)+3$ in the non-symmetric case.


(4) Let $L_h$ be the monomial ideal generated by all monomials of weighted degree
at least $h$.  Then
\[
 I_h=\fkp_H+L_h.
\]
If $h\leq\ell$, every homogeneous element of $\fkp_H$ belongs to $L_h$.
Hence $I_h=L_h$ is a monomial ideal.
Conversely, suppose that $h>\ell$, and choose a nonzero homogeneous binomial in $\fkp_H$ of degree $\ell$.  This binomial belongs to $I_h$, whereas neither of its monomial terms belongs to $I_h$.  Therefore $I_h$ is not a monomial ideal.
\end{proof}


The two types of ideals considered in the main result are both minimally generated by six elements and lie just beyond the monomial case.
In view of Proposition~\ref{prop:inverse-image-ideals} (2) and (3), it is natural to expect that they arise in the following cases.

\begin{itemize}
	\item $H$ has multiplicity three and $h=\ell+1$.
	\item $H$ is symmetric of multiplicity four and $h=\ell+1$.
\end{itemize}

The calculations in the three cases considered in the remainder of this
section show that exactly one defining relation has degree \(\ell\) and that
it remains the unique binomial generator of $I_{\ell+1}$.
For each ideal $I=I_{\ell+1}$ arising in one of these cases, let $f-f'$ denote its unique binomial generator of degree $\ell$.  Since the homogeneous component of $\fkp_H$ of degree $\ell$ is a one-dimensional $k$-vector space spanned by $f-f'$, the monomials $f$ and $f'$ are the only monomials of weighted degree $\ell$.  Passing from $I_{\ell+1}$ to $I_\ell$ amounts precisely to adjoining the degree-$\ell$ monomials $f$ and $f'$.  Since $f-f'\in I_{\ell+1}$, it follows that
$$
I+(f)=I+(f')=I_\ell,
$$
and Proposition~\ref{prop:inverse-image-ideals}(1) shows that this ideal is integrally closed.

We now prove Theorem~\ref{thm:classification}.  In each case, we first determine the Ap\'ery set, the defining ideal $\fkp_H$, and the least relation degree $\ell$.  As in the proof of Proposition~\ref{prop:inverse-image-ideals}, the least elements of $H$ that are at least $\ell+1$ in the residue classes modulo $\mult(H)$ determine the monomial generators of $I_{\ell+1}$ modulo $\fkp_H$.  Combining these monomials with the defining relations gives a generating set of $I_{\ell+1}$.  For the six-generated cases, we then rewrite the parameters to identify the resulting ideals with the two families in the Introduction and give the converse choices of the semigroup parameters.  After treating these two cases, we consider the non-symmetric multiplicity-four case and show that the resulting ideal has seven minimal generators.


Proposition~\ref{prop:inverse-image-ideals} gives only upper bounds for
$\mu_S(I_h)$.  To determine the exact number of generators in the cases
below, we use the following lemma, which separates the contribution of
$\fkp_H$ from that of $t^hk[t]\cap R$.

\begin{lemma}\label{lem:number-of-generators}
For every $h\in\mathbb N$, there is an exact sequence of $k$-vector spaces
\[
0\longrightarrow
\frac{\fkp_H}{\fkp_H\cap(x_1,x_2,x_3)I_h}
\longrightarrow
\frac{I_h}{(x_1,x_2,x_3)I_h}
\longrightarrow
\frac{t^hk[t]\cap R}{\fkm(t^hk[t]\cap R)}
\longrightarrow0.
\]
Consequently,
\[
\mu_S(I_h)
=\dim_k\frac{\fkp_H}{\fkp_H\cap(x_1,x_2,x_3)I_h}
+\mu_R(t^hk[t]\cap R).
\]
In particular, if
\[
\fkp_H\cap(x_1,x_2,x_3)I_h=(x_1,x_2,x_3)\fkp_H,
\]
then
\[
\mu_S(I_h)=\mu_S(\fkp_H)+\mu_R(t^hk[t]\cap R).
\]
\end{lemma}

\begin{proof}
The restriction of $\varphi_H$ to $I_h$ is surjective onto
$t^hk[t]\cap R$ with kernel $\fkp_H$, and
\[
 \varphi_H\bigl((x_1,x_2,x_3)I_h\bigr)
 =\fkm(t^hk[t]\cap R).
\]
Hence $\varphi_H$ induces the asserted exact sequence, and the formulas for
$\mu_S(I_h)$ follow by taking $k$-dimensions.
\end{proof}

In the applications below, the minimal homogeneous generators
$\rho_1,\rho_2,\ldots,\rho_s$ of $\fkp_H$ have pairwise distinct weighted
degrees.  Since $\fkp_H\cap(x_1,x_2,x_3)I_h$ is homogeneous, to verify the
equality
\[
\fkp_H\cap(x_1,x_2,x_3)I_h=(x_1,x_2,x_3)\fkp_H
\]
it is therefore enough to show that
$\rho_i\notin(x_1,x_2,x_3)I_h$ for every $1\leq i\leq s$.



\subsection{Multiplicity three}
\label{subsec:multiplicity-three}
Let
\[
 H=\langle 3,3p+1,3q+2\rangle,\qquad p,q\geq1.
\]
Every numerical semigroup of multiplicity three and embedding dimension
three can be written in this form: after reordering, the two minimal
generators other than $3$ lie in the two nonzero residue classes modulo $3$.
In either of the following membership tests, the coefficient of the generator
not divisible by $3$ must be congruent to $2$ modulo $3$, and its least
possible value is therefore $2$.  It follows that
$3q+2\in\langle3,3p+1\rangle$ if and only if $q\geq2p$, whereas
$3p+1\in\langle3,3q+2\rangle$ if and only if $p\geq2q+1$.  Hence the
integers $3,3p+1,3q+2$ form a minimal system of generators precisely when
\[
 p\leq q\leq2p-1
 \qquad\text{or}\qquad
 q<p\leq2q.
\]
In this case,
\[
 \Ap(H,3)=\{0,3p+1,3q+2\}.
\]

Give $x_1,x_2,x_3$ the weighted degrees $3$, $3p+1$, and $3q+2$, respectively.
The defining ideal is
\[
 \fkp_H=\bigl(x_1^{p+q+1}-x_2x_3,\ x_2^2-x_1^{2p-q}x_3,
                    x_3^2-x_1^{2q-p+1}x_2\bigr).
\]
The degrees of these three relations are
\[
 3(p+q+1),\qquad 6p+2,\qquad 6q+4.
\]
Let $\ell$ be the least of these degrees.


Suppose first that $p\leq q \le 2p-1$.  
Then $\ell = 6p+2$.
Since $\Fr(H)=3q-1 < 6p+3 = \ell + 1$, 
every integer at least $\ell + 1$ belongs to $H$.  Thus the least elements of $H$
at least $\ell + 1$ in the three residue classes modulo $3$ are
\[
 6p+3,\qquad 6p+4,\qquad 6p+5,
\]
represented by
\[
 x_1^{2p+1},\qquad x_1^{p+1}x_2,\qquad x_1^{2p-q+1}x_3,
\]
respectively.  Hence the relation of degree $6p+2$ remains binomial, while
the other two relations contribute only monomial generators to $I_{\ell+1}$, and we
obtain
\[
 I_{\ell+1}=\bigl(x_2^2-x_1^{2p-q}x_3,\ x_1^{2p+1},\ x_1^{p+1}x_2,\
             x_1^{2p-q+1}x_3,\ x_2x_3,\ x_3^2\bigr).
\]

Suppose that $q<p \le 2q$. Then the least relation degree is $\ell = 6q+4$.
In this case $\Fr(H)=3p-2<6q+5 = \ell + 1$, and the same calculation gives
\[
 I_{\ell + 1} =\bigl(x_3^2-x_1^{2q-p+1}x_2,\ x_1^{2q+2},\ x_1^{q+1}x_3,\
             x_1^{2q-p+2}x_2,\ x_2x_3,\ x_2^2\bigr).
\]
In either case, the three monomials obtained from the residue classes have
degrees $\ell+1,\ell+2,\ell+3$ and minimally generate
$t^{\ell+1}k[t]\cap R$.  The three generators of $\fkp_H$ given above have
distinct degrees.  In the first case, the relation of degree $\ell$
does not belong to $(x_1,x_2,x_3)I_{\ell+1}$.  If $q=p$, the other two
relations have degrees $\ell+1$ and $\ell+2$, so neither belongs to this
ideal.  If $q>p$, then, modulo $(x_1,x_2,x_3)I_{\ell+1}$, they are congruent
to $-x_2x_3$ and $x_3^2$, respectively.  Setting $x_1=0$, and then setting
$x_1=x_2=0$, shows that these monomials do not belong to
$(x_1,x_2,x_3)I_{\ell+1}$.  In the second case, the relation of degree
$\ell$ again does not belong to this ideal.  The other two relations are
congruent to $-x_2x_3$ and $x_2^2$, respectively, except that when $p=q+1$
the first of them has degree $\ell+2$.  The same substitutions, with $x_2$
and $x_3$ interchanged for the second monomial, give the required
noncontainments.  Hence
\[
 \fkp_H\cap(x_1,x_2,x_3)I_{\ell+1}=(x_1,x_2,x_3)\fkp_H,
\]
and Lemma~\ref{lem:number-of-generators} gives
\[
 \mu_S(I_{\ell+1})=3+3=6.
\]

To identify these expressions with \eqref{eq:family-one}, we make
the following final choices of variables: $(x,y,z)=(x_1,x_2,x_3)$ in the
first case and $(x,y,z)=(x_1,x_3,x_2)$ in the second case.  With these
choices, both cases are instances of the family
\[
 I=\bigl(y^2-x^{b-d}z,\ x^{2b+2-e},\ x^{b+1}y,\
             x^{b-d+1}z,\ yz,\ z^2\bigr),
 \qquad b>d\geq0,\quad e\in\{0,1\},
\]
where $(b,d,e)=(p,q-p,1)$ for the first case and $(b,d,e)=(q,p-q-1,0)$ for the second case.

Conversely, let $b>d\geq0$ and $e\in\{0,1\}$.  If $e=1$, take
$(p,q)=(b,b+d)$, and if $e=0$, take $(p,q)=(b+d+1,b)$.  In either case,
the required inequalities for $p$ and $q$ hold, and the preceding calculation yields the
ideal in \eqref{eq:family-one}.

The grading inherited from $H$ is given by
$\deg x = 3$, $\deg y = 3b+2-e$, $\deg z = 3b+3d+4-2e$, and hence
$\ell=6b+4-2e$.
These are the weights arising from the numerical-semigroup presentation; the same ideal may admit other positive gradings.

\subsection{Symmetric multiplicity four}
\label{subsec:symmetric-multiplicity-four}
By Corollary~\ref{cor:emb4symm}, $H$ has the form
\[
 H=\langle4,2p,p+2q\rangle,
 \qquad p\geq3\text{ odd},\quad q\geq1.
\]
Its Ap\'ery set with respect to $4$ is
\[
 \Ap(H,4)=\{0,2p,p+2q,3p+2q\}.
\]
Give $x_1,x_2,x_3$ the degrees $4$, $2p$, and $p+2q$, respectively.  The
defining ideal is
\[
 \fkp_H=(x_2^2-x_1^p,\ x_3^2-x_1^qx_2).
\]
The relation degrees are $4p$ and $2p+4q$, so we put $\ell = \min\{4p, 2p+4q\}$.

Assume first that $2q<p$.  Then $\ell=2p+4q$.  Let $w$ be the
integer determined by
\[
 4(w-1)<p+2q+1\leq4w.
\]
The Ap\'ery set of $H$ shows that the least elements at least $\ell + 1$ in
the four residue classes modulo $4$ are represented by
\[
\begin{array}{c|c}
 \text{least element at least }\ell + 1&\text{corresponding monomial}\\ \hline
 2p+4q+2&x_1^{q+\frac{p+1}{2}}\\
 2p+4q+4&x_1^{q+1}x_2\\
 p+2q+4w&x_1^wx_3\\
 3p+2q&x_2x_3.
\end{array}
\]
The last entry is at least $\ell+1=2p+4q+1$ because $2q<p$.
Moreover, since $p$ is odd and $2q<p$, we have
$q+\frac{p+1}{2}\leq p$.
Since
$x_1^p$ is a multiple of $x_1^{q+\frac{p+1}{2}}$, the first relation
of $\fkp_H$ can be replaced by the monomial
$x_2^2$.  We therefore obtain
\[
 I_{\ell+1} =\bigl(x_3^2-x_1^qx_2,\ x_1^{q+\frac{p+1}{2}},\ x_1^wx_3,\ x_1^{q+1}x_2,\ 
             x_2x_3,\ x_2^2\bigr).
\]
To compare this ideal with \eqref{eq:family-one}, we need to rewrite the
exponents in the above generators of $I_{\ell+1}$.
Since $p+2q+1$ is even, exactly one of
$p+2q+1$ and $p+2q+3$ is divisible by $4$.  
Let $e \in \{0,1\}$ such that
\[
 p+2q+1+2e\equiv0\pmod4,
\]
and put
\[
 w=\frac{p+2q+1+2e}{4},\qquad
 b=w-1,\qquad d=w-q-1.
\]
Then $2w-e=q+\frac{p+1}{2}$ and $b>d\geq0$.
Finally, relabelling $(x_1,x_2,x_3)$ as $(x,z,y)$, the ideal becomes the
one in \eqref{eq:family-one}.
Under this relabelling, the grading inherited from $H$ is
$\deg x=4$, $\deg y=p+2q$, and $\deg z=2p$, and the least relation degree is
$\ell=2p+4q$.

Now suppose that $p<2q$.  In this case $\ell =4p$.
If $2q\geq3p+1$, then
$p+2q\geq4p+1=\ell+1$. 
Hence $x_3\in I_{\ell+1}$.  
Put $J=I_{\ell+1}\cap k[x_1,x_2]$.  Then
$I_{\ell+1}=(x_3)+JS$, and $J$ is an integrally closed
$(x_1,x_2)$-primary ideal because it is the contraction of the integrally
closed ideal $I_{\ell+1}$.
In this case, by \cite[Theorem 2.1]{EndoGotoHongUlrich2026}, the Rees algebra of $I_{\ell+1}S_{(x_1,x_2,x_3)}$ is a Cohen--Macaulay normal domain.  Since $I_{\ell+1}S_\fkp=S_\fkp$ for every prime ideal $\fkp\neq(x_1,x_2,x_3)$, the same holds for the Rees algebra of $I_{\ell+1}$.
For this reason, we exclude this case from the family considered in Theorem A. 
We therefore consider the remaining case $p<2q\leq3p$.
Let $r$ be determined by
\[
 4(r-1)+p+2q<4p+1\leq4r+p+2q.
\]
The least elements of $H$ at least $\ell+1$ in the four residue classes modulo $4$ are as follows:
\[
\begin{array}{c|c}
 \text{least element at least }\ell + 1&\text{corresponding monomial}\\ \hline
 4p+4&x_1^{p+1}\\
 4p+2&x_1^{\frac{p+1}{2}}x_2\\
 p+2q+4r&x_1^rx_3\\
 3p+2q&x_2x_3.
\end{array}
\]
Since $2q>p$, the monomial $x_1^qx_2$ is a multiple of
$x_1^{\frac{p+1}{2}}x_2$, so the second relation in
\(\fkp_H\) is replaced by $x_3^2$.  Writing
$p=2b+1$, the inequalities $p < 2q \le 3p$ and
the defining inequality for $r$ give $b\geq1$ and $1\leq r\leq b+1$.
Consequently,
\[
 I_{\ell+1} =\bigl(x_2^2-x_1^{2b+1},\ x_1^{2b+2},\ x_1^{b+1}x_2,\
             x_1^rx_3,\ x_2x_3,\ x_3^2\bigr),
 \qquad b\geq1,\quad1\leq r\leq b+1.
\]

We next verify that the two generating sets obtained above are
minimal.  In each case, the four monomials in the corresponding table
minimally generate $t^{\ell+1}k[t]\cap R$.  Indeed, the degrees of the first
three differ by less than $4$.  A positive difference between the degree of
$x_2x_3$ and either of the first two degrees is an odd integer smaller than
$p+2q$, while its difference from the remaining degree is smaller than $2p$
and congruent to $2$ modulo $4$.  None of these differences belongs to $H$.

The two generators of $\fkp_H$ have distinct degrees.  If $2q<p$, the
relation of degree $\ell$ does not belong to
$(x_1,x_2,x_3)I_{\ell+1}$.  For the other relation, if
$p=2q+1$, its degree is $\ell+2$; otherwise, modulo
$(x_1,x_2,x_3)I_{\ell+1}$, it is congruent to $x_2^2$, which does not belong
to that ideal, as is seen by setting $x_1=x_3=0$.  If $p<2q\leq3p$, the
first relation has degree $\ell$ and the same argument applies to the second:
if $2q=p+1$, its degree is $\ell+2$; otherwise it is congruent modulo
$(x_1,x_2,x_3)I_{\ell+1}$ to $x_3^2$, which does not belong to that ideal,
as is seen by setting $x_1=x_2=0$.
Thus, in both cases,
\[
 \fkp_H\cap(x_1,x_2,x_3)I_{\ell+1}=(x_1,x_2,x_3)\fkp_H.
\]
Lemma~\ref{lem:number-of-generators} now gives
\[
 \mu_S(I_{\ell+1})=2+4=6.
\]
Finally, relabelling $(x_1,x_2,x_3)$ as $(x,y,z)$ gives
an ideal of the form \eqref{eq:family-two}.
The grading inherited from $H$ is given by
$\deg x=4$, $\deg y=2p=4b+2$, and
$\deg z=p+2q$, and the least relation degree is $\ell=4p=8b+4$.
The defining inequality for $r$ shows that $\deg z$ is one of
$8b+5-4r$ and $8b+7-4r$.  The ideal itself is homogeneous for any positive
choice of $\deg z$.

Conversely, let $b\geq1$ and $1\leq r\leq b+1$.  Set
$p=2b+1$ and $q=3b+3-2r$.  Then $p$ is odd, $q\geq1$, and
$p<2q\leq3p$.  Moreover,
\[
 4(r-1)+p+2q<4p+1\leq4r+p+2q,
\]
so the preceding calculation yields the ideal corresponding to the given
pair $(b,r)$.
In Subsection~\ref{sec:family-two}, we use the grading obtained from this
realization; thus $\deg x=4$, $\deg y=4b+2$, and
$\deg z=8b+7-4r$.

For convenience, the parametrizations obtained in the multiplicity-three
and symmetric multiplicity-four cases are summarized below.  
\[
\begin{array}{c|c|c}
 \text{Semigroup}&\text{Range}&\text{Result}\\ \hline\hline
 \langle3,3p+1,3q+2\rangle
   &p\leq q\leq2p-1
   &\eqref{eq:family-one},\quad(b,d,e)=(p,q-p,1)\ \\ \hline
 \langle3,3p+1,3q+2\rangle
   &q<p\leq2q
   &\eqref{eq:family-one},\quad(b,d,e)=(q,p-q-1,0)\ \\ \hline
 \langle4,2p,p+2q\rangle
   &2q<p
   &\eqref{eq:family-one},\quad(b,d,e)=(w-1,w-q-1,e)\ \\ \hline
 \langle4,2p,p+2q\rangle
   &p<2q\leq3p
   &\eqref{eq:family-two},\quad (b,r)=\bigl((p-1)/2,\ \lceil(3p+1-2q)/4\rceil\bigr) \\ \hline
 \langle4,2p,p+2q\rangle
   &2q\geq3p+1
   &I_{\ell+1}=(x_3)+JS
\end{array}
\]

\subsection{Non-symmetric multiplicity four}
\label{subsec:nonsymmetric-multiplicity-four}

Let $H$ be non-symmetric of multiplicity four.  By
Corollary~\ref{cor:emb4symm}, after ordering the generators we may write
\[
 H=\langle4,u,v\rangle,
 \qquad 4<u<v,
\]
where $u$ and $v$ are odd.  The minimality of the generators implies that
$u\not\equiv v\pmod4$ and $v<3u$.  Hence $u+v\equiv0\pmod4$.

Give $x_1,x_2,x_3$ the weighted degrees $4,u,v$, respectively.  By
Theorem~\ref{thm:herzog},
\[
 \fkp_H=\left(
 x_1^{\frac{u+v}{4}}-x_2x_3,
 x_2^3-x_1^{\frac{3u-v}{4}}x_3,
 x_3^2-x_1^{\frac{v-u}{2}}x_2^2
\right).
\]
Moreover,
\[
 \Ap(H,4)=\{0,u,v,2u\}.
\]
The degrees of the three relations are $u+v$, $3u$, and $2v$.  If
$v<2u$, then $u+v$ is the least of these degrees.  If $2u<v<3u$, then
$3u$ is the least.  We therefore treat these two cases separately.

Suppose first that $v<2u$, so that $\ell=u+v$.  Put
\[
 b=\frac{u+v}{4},
\]
and let $e=0$ if $(u,v)\equiv(1,3)\pmod4$ and $e=1$ if
$(u,v)\equiv(3,1)\pmod4$.  Put
\[
 c=\frac{v+2e+1}{4}.
\]
Then
\[
 u=4(b-c)+2e+1,
 \qquad v=4c-2e-1,
\]
and the inequalities $4<u<v<2u$ are equivalent to
\[
 c\geq e+2, \qquad \text{ and }
 \qquad 3c-2e\leq2b\leq4c-2e-2.
\]
Since $\ell+1=4b+1$, the least elements of $H$ that are at least
$\ell+1$ in the four residue classes modulo $4$ are represented as follows:
\[
\begin{array}{c|c|c}
 \text{Ap\'ery element}&\text{least element at least }\ell+1
   &\text{corresponding monomial}\\ \hline
 0&4b+4&x_1^{b+1}\\
 u&4b+2e+1&x_1^cx_2\\
 v&4b-2e+3&x_1^{b-c+1}x_3\\
 2u&4b+2&x_1^{2c-b-e}x_2^2.
\end{array}
\]
The first relation has degree $\ell$ and remains binomial.  Furthermore,
\[
 \frac{3u-v}{4}-(b-c+1)=2b-3c+2e\geq0
\]
and
\[
 \frac{v-u}{2}-(2c-b-e)=2c-b-e-1\geq0.
\]
Thus $x_1^{\frac{3u-v}{4}}x_3$, which appears as the latter monomial term in the second relation, is a multiple
of $x_1^{b-c+1}x_3$, and $x_1^{\frac{v-u}{2}}x_2^2$ is a multiple of $x_1^{2c-b-e}x_2^2$.  Consequently,
\[
 I_{\ell+1}=\bigl(
 x_1^b-x_2x_3,
 x_1^{b+1},
 x_1^cx_2,
 x_1^{2c-b-e}x_2^2,
 x_1^{b-c+1}x_3,
 x_2^3,
 x_3^2
 \bigr).
\]

Suppose next that $2u<v<3u$, so that $\ell=3u$.  Let $e=0$ if
$u\equiv1\pmod4$ and $e=1$ if $u\equiv3\pmod4$, and put
\[
 c=\frac{u-2e-1}{4},
 \qquad b=\frac{3u-v}{4}.
\]
Then
\[
 u=4c+2e+1,
 \qquad v=3u-4b,
\]
and the inequalities $2u<v<3u$ are equivalent to
\[
 c\geq1,
 \qquad 1\leq b\leq c.
\]
Since $\ell+1=3u+1$, the least elements in the four residue classes are
represented as follows:
\[
\begin{array}{c|c|c}
 \text{Ap\'ery element}&\text{least element at least }\ell+1
   &\text{corresponding monomial}\\ \hline
 0&\ell+1+2e&x_1^{3c+2e+1}\\
 u&\ell+2&x_1^{2c+e+1}x_2\\
 2u&\ell+3-2e&x_1^{c+1}x_2^2\\
 v&\ell+4&x_1^{b+1}x_3.
\end{array}
\]
The second relation has degree $\ell$ and remains binomial.  The equalities
\[
 \frac{u+v}{4}-(3c+2e+1)=c-b\geq0
\]
and
\[
 \frac{v-u}{2}-(c+1)=3c+2e-2b\geq0
\]
show that the other two relations may be replaced by $x_2x_3$ and $x_3^2$.
We therefore obtain
\[
 I_{\ell+1}=\bigl(
 x_2^3-x_1^bx_3,
 x_1^{3c+2e+1},
 x_1^{2c+e+1}x_2,
 x_1^{c+1}x_2^2,
 x_1^{b+1}x_3,
 x_2x_3,
 x_3^2
 \bigr).
\]

In both cases, the four monomials obtained from the four residue classes have
weighted degrees $\ell+1,\ell+2,\ell+3,\ell+4$, in some order.  Their images
therefore form a minimal system of generators of
$t^{\ell+1}k[t]\cap R$, since the least positive element of $H$ is $4$.
The generating sets of $I_{\ell+1}$ obtained above also satisfy
\[
 \fkp_H\cap(x_1,x_2,x_3)I_{\ell+1}
 =(x_1,x_2,x_3)\fkp_H.
\]
Indeed, the relation of degree $\ell$ cannot belong to
$(x_1,x_2,x_3)I_{\ell+1}$, whose nonzero homogeneous elements have degree at
least $\ell+4$.  In the case $v<2u$, if
either $2b-3c+2e$ or $2c-b-e-1$ is positive, the corresponding relation is
congruent modulo $(x_1,x_2,x_3)I_{\ell+1}$ to $x_2^3$ or $x_3^2$,
respectively.  If the difference is zero, that relation has degree at most
$\ell+3$.  The substitutions
$x_1=x_3=0$ and $x_1=x_2=0$ show that neither $x_2^3$ nor $x_3^2$ belongs to
$(x_1,x_2,x_3)I_{\ell+1}$.  In the case $2u<v<3u$, the third relation is
congruent to $x_3^2$.  If $c>b$, the first relation is congruent to
$-x_2x_3$; if $c=b$, it has degree at most $\ell+3$.  After setting $x_1=0$,
the monomial $x_2x_3$ does not belong to
$(x_2,x_3)(x_2^3,x_2x_3,x_3^2)$, and the assertion for $x_3^2$ follows by
setting $x_1=x_2=0$.  Since the three relation degrees are distinct,
this proves the claim.
Lemma~\ref{lem:number-of-generators} now gives
\[
 \mu_S(I_{\ell+1})=3+4=7.
\]
This proves Theorem~\ref{thm:classification}(3), and the converse constructions in the preceding two
subsections prove Theorem~\ref{thm:classification}(4).  The proof of Theorem~\ref{thm:classification} is complete.
